\documentclass[11pt]{article}
\usepackage[a4paper,margin=25mm]{geometry}
\usepackage[T1]{fontenc}
\usepackage{lmodern}
\usepackage{booktabs}
\usepackage{hyperref}
\usepackage{xcolor}
\usepackage{listings}
\lstset{basicstyle=\ttfamily\small, columns=fullflexible, keepspaces=true,
  breaklines=true, frame=single, framerule=0.3pt, rulecolor=\color{gray}}
\newcommand{\code}[1]{\texttt{#1}}
\newcommand{\us}{\textmu s}

\title{\code{wfview.py}: a viewer for the summed waveforms}
\author{BGOTSCAN}
\date{2026-10-14}

\begin{document}
\maketitle

\section{Purpose}

A first look at the data, event by event: the sum of the 64 SiPM waveforms,
raw and baseline-subtracted, with an optional moving average; the individual
channels of one event; and the average summed pulse of a run. It reads the
raw HDF5 files with h5py (the file layout is described in
\code{jscan\_doc.tex}) and does no analysis beyond a baseline and a moving
average. It is a Python script in \code{src/py/}.

\section{Usage}

\begin{lstlisting}
python3 src/py/wfview.py RUN [options]
\end{lstlisting}

\code{RUN} is a run name (\code{128K\_37V\_th19}), with or without
\code{.h5}, looked up in the data directory, or a path to a file.

\begin{center}
\begin{tabular}{l l p{80mm}}
\toprule
option & default & meaning \\
\midrule
\code{--data DIR} & \code{2026\_10\_02\_bgo\_cryo\_tempscan/data\_cryo} (relative to the repository) & directory of the run files \\
\code{--events SPEC} & \code{0-3} & events to show: a range \code{a-b} (inclusive) or a list \code{a,b,c}; 0-based as in the file \\
\code{--baseline-window N} & 150 & number of pre-trigger samples, from the start of the record, used for the baseline \\
\code{--smooth N} & 0 & width of the moving average in samples; 0 = off \\
\code{--channels} & off & also show the 64 channels of the first selected event \\
\code{--average} & off & also show the average summed pulse of the whole run \\
\code{--save NAME} & off & save the figures as PNG under \code{src/py/plots/RUN/} with the prefix \code{NAME}; without it the figures are shown on screen \\
\bottomrule
\end{tabular}
\end{center}

\section{Definitions}

\paragraph{Time axis.} $t_i = (i - \code{pretrigger})\,\code{tsample\_ns}$,
with $i$ the 0-based sample index, so the trigger is at $t = 0$ and the record
runs from $-1.6$~\us{} to $22.4$~\us{} for these runs.

\paragraph{Baseline.} Two baselines are used.
\begin{itemize}
\item \emph{WebDump}: the per-channel values stored in the attribute
  \code{baselines}, measured with software triggers at the start of the run.
\item \emph{Pre-trigger}: per event and channel, the median of the first $N$
  samples of the record (\code{--baseline-window}, default 150 of the 200
  pre-trigger samples). The window stops before the trigger because the
  pulse often starts before $t = 0$ (the trigger fires late on these slow,
  small pulses); keep $N$ small enough that the window is clean, and check
  it on the average pulse.
\end{itemize}

\paragraph{Summed waveform.} $S_i = \sum_c (w_{c,i} - b_c)$ over the 64
channels, with $b_c$ one of the baselines above. With the pre-trigger
baseline, the pre-trigger level of $S$ is zero by construction, event by
event.

\paragraph{Moving average.} With \code{--smooth N}, $S$ is replaced by its
mean over a window of $N$ samples centred on each sample (edges use the
samples available). The white noise of the sum (about 30 ADC counts rms for
64 channels at 3.9 counts each) goes down by $\sqrt{N}$; a pulse of several
\us{} is not distorted by $N = 64$ (0.5~\us), a 300~ns pulse would be.

\section{Figures}

\paragraph{Summed waveforms (always).} One panel per selected event, in a
grid. Each panel shows $S$ with the WebDump baselines (thin, grey), $S$ with
the pre-trigger baseline (blue) and, if requested, its moving average
(orange). The baseline window is shaded, the trigger is a dashed line. The
title gives the event number and, from the pre-trigger-baseline sum: the
maximum, its time, and the integral over the whole record in ADC\,$\cdot$\,ns.
The same numbers are printed on the terminal, one line per event.

\paragraph{Channels (\code{--channels}).} The 64 channels of the first
selected event, baseline-subtracted (pre-trigger), on an $8 \times 8$ grid
in channel order, shared axes, with the channel number in each panel.

\paragraph{Average pulse (\code{--average}).} The mean of $S$ over all events
of the run, read in chunks, with both baselines; the pre-trigger-baseline
version is the reference. This is the figure that shows the pulse shape at
each temperature and tells whether the baseline window is clean.

\section{Examples}

\begin{lstlisting}
# four events at 128 K, smoothed over 0.5 us
python3 src/py/wfview.py 128K_37V_th19 --events 0-3 --smooth 64

# one event with its 64 channels
python3 src/py/wfview.py 128K_37V_th19 --events 7 --channels

# average pulse of the warm run, saved to src/py/plots/187K_38d5V_th19/187K_*.png
python3 src/py/wfview.py 187K_38d5V_th19 --average --save 187K
\end{lstlisting}

\end{document}
